\documentclass[11pt,a4paper]{article}
\usepackage[T1]{fontenc}
\usepackage[utf8]{inputenc}
\usepackage[margin=25mm]{geometry}
\setlength{\emergencystretch}{3em}
\usepackage{amsmath,amssymb,booktabs,tabularx,microtype}
\usepackage[colorlinks=true,allcolors=blue]{hyperref}
\title{Calibration and domain shift in four-mode infrared focal-plane wavefront sensing}
\author{Mojtaba Taheri}
\date{Working draft 0.4 --- 2 October 2026}
\newcommand{\nm}{\,\mathrm{nm}}
\begin{document}
\maketitle
\noindent\textbf{Draft status.} Baseline results below are verified local computational
results, including the bounded classical comparison, fresh detector challenges
completed independent-engine validation, matched architecture controls and
paired shift-factor diagnostics. Untouched confirmation remains unfinished.
This document is not ready for submission. Affiliation, authorship contributions,
funding, conflicts and venue-specific disclosures require author review.

\begin{abstract}
Fast inference alone does not establish reliable infrared low-order wavefront
sensing. We study signed focus, two astigmatism components and spherical
aberration in an integrated broadband, undersampled focal-plane benchmark.
The experiment uses a published Keck pupil proxy and a filtered atmospheric
residual surrogate, with independent atmospheric and calibration parents across
training, tuning, interval calibration and evaluation. Frozen two-layer neural
estimators have nominal four-mode wavefront errors of 60.38 and 63.05 nm for
single and paired measurement designs, compared with 80.26 and 76.95 nm for
ridge regression. A combined residual-spectrum, amplitude and diversity shift
raises neural errors to approximately 106 nm; paired standardized-MSE contrasts
with ridge become inconclusive. Empirical interval coverage also declines.
On a smaller 10,000-electron subset, nonlinear physical fits improve nominal
error at much greater computational cost, but physical refinement worsens
shifted error. Matched ResNet18 controls reduce nominal errors to 53.68 and
57.17 nm, at approximately nine times the MLP median inference cost, without
a consistent shifted advantage. A twenty-parent paired factorial study
identifies residual amplitude and spectrum as the dominant average degradation
factors within its tested ranges. Diversity effects are smaller and uncertain
for the neural estimators. These development results establish neither
telescope performance nor superiority over the author's prior pipeline;
independently seeded confirmation remains necessary.
\end{abstract}

\section{Motivation and relation to previous work}
Low-order aberrations affect image position, focus and point-spread-function
structure, limiting the stability needed for high-contrast observations.
Measuring errors near the science focal plane reduces the separation between
the sensed optical path and the path whose image quality matters. Phase
retrieval and phase diversity encode aberration information through intensity
and known optical changes \cite{fienup,paxman}. Photon statistics, detector
sampling and calibration constrain the information available to any estimator
\cite{guyon}. LIFT provides a linearized route to low-order estimation
\cite{meimon,plantet2011,plantet2013}. Sequential techniques such as Fast \&
Furious use an image history and control diversity \cite{korkiakoski,bos}.
These measurement resources must be included when comparing estimators.
Photon noise, sampling and chromatic effects constrain useful sensing
performance independently of reconstruction speed \cite{guyon}.

Infrared sensing is relevant for cool or dust-obscured targets that are faint
in the visible but sufficiently bright in the infrared. Keck II's infrared
pyramid demonstrates this motivation \cite{bond}. The benefit depends on target
spectrum, throughput, detector properties and residual aberrations; infrared
operation alone does not establish sensitivity or sky coverage for this benchmark.

LOWFS includes distinct tasks: slow sodium-focus truth sensing, fast
tip--tilt--focus control, non-common-path aberration calibration, and differential
segment or petal phasing. Their information, cadence and validation requirements
differ. The present experiment estimates static signed focus, two astigmatism
components and spherical aberration from integrated images. It does not validate
tip--tilt control, petal phasing or a telescope control loop.

The closest infrared learned precedent is Taheri et al.\ \cite{taheri}, which
uses ResNet18 with H-band simulations and K-band Keck internal calibration-source
measurements. Those measurements are not atmospheric on-sky validation.
The published cross-band focus comparison removes a mean offset, and therefore
does not by itself establish absolute focus calibration.
Salgueiro et al.\ \cite{salgueiro} subsequently selected Gerchberg--Saxton for
2025 Keck slow-focus on-sky tests because of its residual robustness in their
comparison. The paper describes further operational assessment as future work.
Kuznetsov et al.\ \cite{kuznetsov} show that independent NCPA calibration can
substantially improve LIFT focus estimation in VLT/MUSE-NFM tests. These studies
make calibration and residual mismatch central comparison variables.

Convolutional focal-plane reconstruction, engineered phase diversity and
accelerated phase-diversity networks already have substantial precedent
\cite{orban,quesnel,wu}. Residual networks supply a general architecture
framework \cite{he}; they do not guarantee an advantage for undersampled
infrared images. Learned reconstruction has also reached on-sky operation
with an unmodulated pyramid \cite{landman}; preliminary learned infrared
Lyot-based LOWFS has controlled differential piston on SCExAO \cite{singh}.
These instruments encode phase through additional optics and receive different
measurements from a direct focal-image estimator. They establish relevant
precedent without a photon-equivalent algorithm ranking. Reconstruction timing
must also be distinguished from exposure-to-actuation latency.

The question here is whether additional spatial structure or
physical refinement improves error, uncertainty and measured cost under a
matched measurement and calibration protocol. The completed baseline is a
development result. New methods selected after inspecting it require fresh,
independently seeded confirmation.

\section{Measurement and computational experiment}
\subsection{Targets and pupil convention}
The target vector is
$\boldsymbol a=(a_f,a_{a_c},a_{a_s},a_s)$: focus, astigmatism cosine,
astigmatism sine and spherical aberration. Coefficients are signed nm OPD RMS
on a frozen pupil-weighted basis. Conventional polynomial seeds, after removal
of their pupil means, are orthonormalized jointly with physical tip and tilt
using a stored Cholesky transformation. The reference metric is normalized
geometric illuminated area. Global piston is excluded. Four-mode error is
distinct from the remaining high-order AO wavefront error.

With normalized pupil modes $q_k$, known diversity $d_c$, calibrated static
aberration $W_0$ and unknown disturbance $r$, the pupil field is proportional to
\begin{equation}
 A(\boldsymbol x)\exp\!\left[\frac{2\pi i}{\lambda}
 \left(10^{-9}\sum_{k=1}^4 a_kq_k(\boldsymbol x)
             +W_0(\boldsymbol x)+r(\boldsymbol x,t)+d_c(\boldsymbol x)\right)\right],
\end{equation}
where all OPD terms inside the brackets are in metres. Detector expectations
integrate the propagated intensity over wavelength, detector pixel area and
exposure time. Full-plane normalization is retained; the detector crop is not
renormalized to unit flux. Inference receives count images and frozen
calibration, with no target coefficients, centroid truth, atmospheric phase
realization or per-parent calibration-error map.

\subsection{Historical Keck proxy and acquisition}
The optical generator uses HCIPy 0.7.1 \cite{por} and its published segmented
Keck pupil, physically scaled to 10.949 m. The library's physical gap setting is
used with geometric edge quadrature. This is a published pupil proxy, rather
than an as-built Keck I/TRICK calibration. A 16$\times$16 region has 50 mas pixels;
the detected-photon band is 1.5--1.8$\,\mu$m. Five wavelength nodes and
six-point quadrature on each pixel axis are used at the admitted sampling.
The 1024-square pupil, wavelength and pixel sampling were refined jointly
before generation. These are modeled choices, not current observatory specifications.

One measurement design uses a 10 ms exposure with 200 nm astigmatism diversity.
The other uses two sequential 5 ms sub-exposures, first focused and then with
200 nm focus diversity. Both have the same total source-electron budget
$F\in\{10^3,10^4,10^5\}$. Each paired image receives half that budget and half
the 5.01-electron background-plus-dark expectation per pixel, but incurs its
own 5-electron RMS read contribution. Signed read-noise counts are retained.
The two designs therefore share exposure time and total photons, while their
temporal sampling, diversity and number of reads differ.

\subsection{Residual and calibration model}
Seven frozen-flow layers use the atmosphere parameters reported by
Taheri et al.\ \cite{taheri}: $r_0=0.165$ m, $L_0=75$ m and the published layer
heights, fractions and winds. The fractions sum to 0.97 and are explicitly
renormalized. A 500 nm reference for $r_0$ is a modeling assumption. A 128-square
atmospheric grid is interpolated to the optical pupil. The declared spatial
residual amplitude filter is
\begin{equation}
 H(k)=\frac{k^2}{k^2+k_c^2}.
\end{equation}
Independent developer sequences set global gains for 100/200 nm nominal residual
levels. Piston, tilt and all four target modes are projected out of the
disturbance. Sequence and frame fluctuations remain; frames are not individually
RMS-normalized. The same parent supplies twenty consecutive 10 ms acquisitions.
Targets remain constant within an acquisition and vary across acquisitions;
this does not model a realistic closed-loop target time series.
Each acquisition has independent uniform coefficient draws: focus over
$[-150,150]$ nm and each of the other three modes over $[-100,100]$ nm.
Nominal parents are allocated equally to the 100 and 200 nm residual levels.

The inference calibration includes a 50 nm static high-order pattern. Each parent
also has an independently seeded unknown high-order calibration error with RMS
drawn up to 20 nm. Centroids vary within $\pm20$ mas per axis. The shifted evaluation
combines a 300 nm residual level, a cutoff change from $10/D$ to $6/D$, and 10\% excess
diversity. This combined shift cannot identify which component causes the
observed degradation. The residual is a filtered surrogate with no measured
WFS/DM controller, telemetry or scintillation.

\subsection{Grouped data and frozen estimators}
There are 620 independent atmospheric/calibration parents: 300 training,
60 validation, 60 interval-calibration, 100 nominal evaluation and 100 shifted
evaluation. Twenty acquisitions and three flux replicas per parent produce
12,400 physical acquisitions and 37,200 rows per design. The 18,000 training rows
represent 300 independent parents. Replicas and both measurement designs remain
within the same parent split.

Each design has a separate MLP with two 32-unit ReLU hidden layers and four
outputs: 9,412 single-image or 17,604 paired-image parameters. The features are
\begin{equation}
 z_j=\frac{\operatorname{asinh}(y_j/100)-\mu_j}{\sigma_j},
\end{equation}
with means and scales fitted on training parents alone. Outputs are scaled by
$(150,100,100,100)$ nm. Training uses Adam \cite{adam} at learning rate 0.001,
batch 128, standardized coefficient MSE, a 100 epoch maximum and ten stale
validation epochs. Seeds 11, 23, 37 are trained separately. Parent-weighted
validation loss selects seed 11 for single images and 37 for pairs. Ridge
regularization is selected on the same validation set. All model and
preprocessing artifacts are checksummed and frozen before evaluation.

Independent calibration parents define per-mode empirical 95\% absolute-error
quantiles with equal parent weights. These are marginal empirical intervals,
not joint intervals or a finite-sample coverage guarantee. No evaluation
outcome changes their widths in the frozen baseline.

The matched architecture extension uses the same features, label scaling,
splits, Adam settings, batch size, maximum epochs, patience and three seeds.
A compact CNN has a 16-channel convolutional stem, two 16-channel residual
blocks separated by $2\times$ average pooling, and a 32-unit output head:
42,372/42,516 single/pair parameters. ResNet18 has a $3\times3$ stride-one
one/two-channel stem, no initial max-pooling and a four-output head:
11,169,732/11,170,308 parameters. Neither resizes pixels or downloads weights.
Validation selects CNN seeds 37/11 and ResNet18 seeds 11/11 for single/pair.
Each model's widths use the original independent calibration split.
This is a documented adaptation, not an exact reproduction of the 2024 pipeline.

\subsection{Metrics and uncertainty}
For $P$ parents, four-mode reconstructed error is
\begin{equation}
 R=\left[\frac1 P\sum_{p=1}^P\frac1{n_p}
                \sum_{i=1}^{n_p}\|\widehat{\boldsymbol a}_{pi}-\boldsymbol a_{pi}\|_2^2
          \right]^{1/2}.
\end{equation}
Parent means, rather than rows treated as independent, also define coefficient
bias, marginal coverage and paired standardized-MSE contrasts. Reported
standardized MSE is
\begin{equation}
 S=\frac1P\sum_{p=1}^P\frac1{n_p}\sum_{i=1}^{n_p}
       \frac14\sum_{k=1}^4\left(\frac{\widehat a_{pik}-a_{pik}}{s_k}\right)^2,
 \qquad \boldsymbol s=(150,100,100,100)\ \mathrm{nm}.
\end{equation}
This weights modes differently from $R$; a paired confidence interval for
an $S$ contrast does not directly test an unscaled wavefront-RMSE improvement.
For aggregate three-flux results $n_p=60$; for one flux $n_p=20$.
Percentile intervals resample the analysis's parents with replacement for
1,000 draws: 100 per native baseline evaluation split, ten per classical
condition, twenty per detector-challenge condition, 100 per architecture
evaluation split and twenty paired diagnostic parents.
They condition on the frozen fitted models and do not include training-set or
model-selection uncertainty. Flux-specific errors use equal parent weights
within each flux; the equal-third flux aggregate averages squared errors before
taking the square root.

Warmed latency measurements use 1,000 batch-one calls after ten warm-up calls.
CUDA synchronization brackets each call. Preprocessing, device transfer,
inference and output transfer are included; acquisition, detector readout,
diversity switching and controller/actuator latency are excluded.

\section{Verified baseline results}
\begin{table}[ht]
\centering\small
\caption{Four-mode error in nm OPD RMS. Parent-bootstrap intervals condition on
the frozen models. Each evaluation split contains 100 parents.}
\begin{tabular}{@{}llrrr@{}}
\toprule
Condition&Design&MLP&MLP 95\% interval&Ridge\\
\midrule
Nominal&Single&60.38&58.80--61.98&80.26\\
Nominal&Pair&63.05&61.84--64.19&76.95\\
Shifted&Single&106.85&105.12--108.48&109.04\\
Shifted&Pair&106.10&104.43--107.62&107.66\\
\bottomrule
\end{tabular}
\label{tab:baseline}
\end{table}

Nominal paired MLP-minus-ridge standardized-MSE intervals are
$[-0.05389,-0.04914]$ for single images and $[-0.03908,-0.03532]$ for pairs.
Both favor the MLP in standardized MSE on this admitted nominal population. The shifted intervals,
$[-0.00329,0.00404]$ and $[-0.00320,0.00553]$, include zero. These intervals do
not establish a clear advantage under the combined shift.

Nominal single/pair MLP errors at 1,000 electrons are 81.14/93.15 nm;
at 10,000 electrons they are 47.67/45.43 nm; at 100,000 electrons they are 45.61/34.41 nm.
An aggregate ordering therefore does not describe every flux regime. In
particular, the paired design's brighter-flux accuracy must be weighed against
its additional read contribution and different temporal encoding.

\begin{table}[ht]
\centering\small
\caption{Empirical marginal interval coverage; mode order is focus,
astigmatism cosine, astigmatism sine, spherical. Widths remain frozen.}
\begin{tabular}{@{}llrrrr@{}}
\toprule
Condition&Design&Focus&Astig.cos&Astig.sin&Spherical\\
\midrule
Nominal&Single&94.40\%&95.68\%&95.40\%&95.45\%\\
Nominal&Pair&94.97\%&95.35\%&95.05\%&94.12\%\\
Shifted&Single&66.33\%&75.75\%&75.77\%&82.42\%\\
Shifted&Pair&76.05\%&80.60\%&82.98\%&81.38\%\\
\bottomrule
\end{tabular}
\label{tab:coverage}
\end{table}

No nonfinite baseline prediction occurred. This does not imply that every
estimate is scientifically accurate. The faint nominal regime also loses
coverage: approximately 86.7--90.4\% for single images and 83.9--86.3\% for pairs.
Calibration over a flux mixture does not guarantee nominal coverage at every flux.

Nominal single/pair MLP CUDA latency medians on an NVIDIA GeForce RTX 4060 Laptop GPU are 0.341/0.573 ms, with 95th
percentiles 0.848/1.140 ms and 99th percentiles 1.216/2.792 ms. CPU ridge medians
are 0.0118/0.0155 ms. These are hardware-specific implementation measurements;
they do not show that a neural network is required for fast reconstruction.
The frozen six-network training completed in 202.344 s. A supervisor metadata
failure prevented final training CPU/peak-memory recording; a conservative
5,041 CPU-second charge was retained instead of refitting the models.

\section{Bounded classical comparison}
The five-method comparison uses identical 10,000-electron observations from
ten nominal and ten shifted parents, one acquisition per parent. These are
subsets of the existing evaluation data. They support development comparisons,
rather than fresh confirmation of a subsequently chosen method.

\begin{table}[ht]
\centering\small
\caption{Four-mode error in nm for the bounded classical subset. All 200 planned
estimates are finite; no time or memory cap occurred. Each cell has ten parents.
These point errors do not by themselves establish a significant ordering.
Final statuses are OK; one selected coarse capture reached its iteration cap
before successful final refinement.}
\begin{tabular}{@{}lrrrr@{}}
\toprule
Method&Nominal single&Nominal pair&Shifted single&Shifted pair\\
\midrule
MLP&59.09&65.93&73.93&82.99\\
Ridge&76.36&79.13&90.67&93.45\\
LiFT-style&28.12&80.39&133.92&193.05\\
Physical multistart&22.68&21.88&107.14&153.13\\
MLP plus two updates&49.61&45.73&103.34&150.21\\
\bottomrule
\end{tabular}
\end{table}

The calibrated LiFT-style iterative Jacobian fit uses up to five updates and
is not an exact upstream LIFT reproduction. The physical multistart fit uses
three initializations on a 512-square coarse pupil, then refines its best
candidate on the 1024-square pupil with one initialization. Each fit has a
30-function-evaluation cap per initialization. The hybrid uses exactly two physical updates from
the MLP estimate. Every planned observation remains in the denominator;
missing estimates receive a declared penalty of twice each target limit,
with conditional available-estimate errors reported separately.
The selected coarse fit for shifted parent0000 in the single-image design
reached its 30-evaluation cap, while final refinement returned OK. Reported
top-level cap fractions concern final estimates; this nested event is retained
separately. Statuses of discarded coarse starts were not saved.

Nominal single-image LiFT-style and multistart fits have lower standardized
MSE than the MLP: paired parent-bootstrap 95\% contrast intervals are
$[-0.1154,-0.0230]$ and $[-0.1235,-0.0249]$, respectively.
For nominal pairs, the hybrid and multistart intervals are
$[-0.0901,-0.0198]$ and $[-0.1511,-0.0381]$.
Under the combined shift, each physical method has a positive standardized-MSE
contrast interval relative to the MLP in both designs. The physical updates
therefore do not provide a general robustness improvement on this subset.

Observed median call times are approximately 8.7/17.3 s for single/pair
LiFT-style fits, 53.2/99.1 s for multistart and 3.5/6.9 s for the hybrid.
The physical fits use CPU propagation; the hybrid combines a CUDA MLP call
with CPU physical updates.
These implementation timings are from ten calls per condition, interleaved
with other methods. They are not comparable to the separate 1,000-call warmed
tail-latency measurement without this sampling qualification, and do not
establish the best achievable runtime of a physical solver.

\section{Fresh detector and optical challenges}
Twenty independently seeded parents each provide one 10,000-electron acquisition
in nine conditions. Residual amplitude alternates between the nominal 100 and
200 nm tiers. Conditions share each parent's target, centroid, atmospheric
exposure and unknown high-order calibration map; their detector noise draws
are independent. Frozen MLP and ridge estimators receive the same images in
each design and retain their nominal inference calibration. There are 720
planned estimates: 640 finite estimates and 80 expected nonfinite-input rejections,
with no time or memory cap.

\begin{table}[ht]
\centering\small
\caption{Four-mode error in nm for fresh challenges. Each finite cell has twenty
parents. These are point estimates, rather than tests of differences between
conditions. The bad-pixel condition has no estimate for any parent or method;
all twenty planned observations per cell remain in failure denominators.}
\begin{tabular}{@{}lrrrr@{}}
\toprule
Condition&MLP single&MLP pair&Ridge single&Ridge pair\\
\midrule
Nominal&50.37&48.23&63.00&61.82\\
10-electron read noise&57.78&62.31&63.10&61.47\\
Interpixel coupling&51.59&44.32&64.78&58.88\\
2,000-electron clipping&49.53&47.06&62.33&56.52\\
1.1 diversity multiplier&41.97&52.82&57.96&70.16\\
1.02 pixel-scale multiplier&50.14&46.07&61.71&59.72\\
0.01$D$ pupil registration&49.63&50.59&61.15&65.27\\
20 nm extra calibration error&53.03&55.49&64.13&63.65\\
Nonfinite bad pixel&---&---&---&---\\
\bottomrule
\end{tabular}
\label{tab:challenges}
\end{table}

The coupling kernel redistributes two percent of charge to each of four
neighbors before read noise; influx from outside the crop is omitted.
Clipping is a 2,000-electron ceiling on the noisy count image, not a calibrated
physical detector saturation model. Pupil displacement is $0.01D$ along one
axis. The extra calibration pattern has 20 nm projected high-order RMS in
addition to each parent's nominal unknown map. A nonfinite pixel at index
$(3,9)$ invokes the common acquisition guard; no reconstruction or interpolation
is attempted. This tests rejection behavior, not bad-pixel tolerance.
These limited challenges do not certify a characterized detector or establish
robustness to general combinations of errors.

\section{Independent simulator validation and remaining gates}
The admitted native physics checks include signed noiseless recovery,
independent direct propagation, joint pupil/pixel/wavelength refinement,
disturbed exposure refinement and detector count moments. Maximum noiseless
coefficient RMSE is 0.134 nm; admitted joint sampling differences are below
relative image $L_1<10^{-3}$. These checks establish numerical behavior for the
specified forward model, not telescope calibration fidelity.

Full OOPAO validation completed with 100 independent nominal parents, five
physical acquisitions per parent and native exit status zero. The adapter pins source revision
\texttt{e8e9aa60cf99f4ab21a4dae7c29aae9b9ec6ec87}; native propagation and
independently evolved atmospheric screens differ from HCIPy. Pupil/basis,
bandpass, pixel quadrature and detector statistics remain shared experimental
inputs. All seven identical-input optical cases passed, with maximum relative
image $L_1=1.08\times10^{-6}$ and padding-refinement difference
$1.06\times10^{-6}$. Developer radial-PSD $L_1=0.00724$ and lag-one-correlation
difference $0.000918$ passed the declared practical thresholds of 0.35 and 0.05.
These checks do not prove identical stochastic processes or a realistic
telescope controller. The independent simulator has its own fixed developer
gain; no test-based normalization or refitting was used.

The 500 OOPAO acquisitions contain 200, 200 and 100 observations at 1,000,
10,000 and 100,000 electrons. Direct aggregate comparisons would mix fluxes
differently from the original equal-third benchmark. With equal parent weights
within each flux and one-third flux weights, frozen MLP single/pair RMSE is
61.57/63.44 nm and ridge RMSE is 81.38/78.22 nm. MLP errors at the three fluxes
are 84.18/46.67/45.91 nm for single images and 92.12/48.40/35.27 nm for pairs.
No nonfinite prediction occurred. These point summaries support nominal
cross-simulator consistency; they do not quantify an uncomputed paired
contrast or establish robustness under the combined physical shift.
The complete run used 6,698 wall seconds, 6,874.5 charged CPU seconds and
7.16 GiB sampled peak working set on the stated host.

\section{Matched architecture development comparison}
All twelve CNN/ResNet18 fits completed with native exit status zero and frozen
checkpoint, validation-selection and calibration identities. Native admission
included actual CUDA profiles, bitwise checkpoint round trips and separate
interruption/resume. Scientific fitting took 1,804s wall; the full initial-fitting
envelope, including baseline and admission, used 2,078s of its 7,200s ceiling.
The old nominal/shifted test sets are development evidence for this extension.

\begin{table}[ht]
\centering\small
\caption{Validation-selected architecture controls. Four-mode nm OPD RMS,
equal parent and one-third flux weights. OOPAO tests nominal conditions only.}
\begin{tabular}{@{}lrrrrrr@{}}
\toprule
&\multicolumn{2}{c}{Original nominal}&\multicolumn{2}{c}{Original shifted}&\multicolumn{2}{c}{OOPAO}\\
Method&Single&Pair&Single&Pair&Single&Pair\\
\midrule
MLP&60.38&63.05&106.85&106.10&61.57&63.44\\
Ridge&80.26&76.95&109.04&107.66&81.38&78.22\\
Compact CNN&58.52&61.20&108.17&106.24&58.29&62.27\\
ResNet18&53.68&57.17&106.01&107.19&53.74&58.19\\
\bottomrule
\end{tabular}
\label{tab:architectures}
\end{table}

Nominal candidate-minus-MLP conditional 95\% paired wavefront-RMSE intervals
favor CNN: single $[-2.46,-1.27]$ nm, pair $[-2.36,-1.36]$ nm; and ResNet18:
single $[-7.35,-6.11]$ nm, pair $[-6.54,-5.19]$ nm. Standardized-MSE contrasts
also favor both families. Independent OOPAO point results show the same
nominal ordering, without establishing telescope transfer.

Shifted contrasts are mixed. For paired images, CNN improves standardized MSE:
its conditional 95\% interval is $[-0.00958,-0.00139]$. The wavefront contrast
is inconclusive, with interval $[-0.78,1.04]$ nm. ResNet18 pair wavefront RMSE
worsens by 1.09 nm, with interval $[0.18,2.13]$ nm, while its standardized-MSE
interval includes zero. Its single-image wavefront contrast is also
inconclusive. Increasing capacity therefore does not establish robustness
under the combined shift.

Table~\ref{tab:architecture-latency} reports the same warmed 1,000-call
batch-one protocol across methods. The nominal ResNet18 improvement
costs approximately nine times the MLP median inference latency on this host.
These figures exclude acquisition and control and differ from the earlier
baseline timing invocation; the present comparisons share one measurement protocol.

\begin{table}[ht]
\centering\small
\caption{Inference latency in ms on the stated host. Includes preprocessing,
device transfers and synchronization; excludes acquisition and controller.}
\begin{tabular}{@{}lrrrr@{}}
\toprule
Method&Single median&Single p95&Single p99&Pair median\\
\midrule
MLP&0.423&0.623&0.815&0.441\\
Ridge&0.0144&0.0153&0.0213&0.0201\\
Compact CNN&1.012&1.304&2.001&1.021\\
ResNet18&3.871&5.238&7.177&3.928\\
\bottomrule
\end{tabular}
\label{tab:architecture-latency}
\end{table}

\section{Paired residual and diversity diagnostics}
Twenty new development parents cross residual amplitude 200/300nm, spectral
cutoff$10/D$/$6/D$ and diversity multiplier1.0/1.1: eight cells per parent.
Targets, centroid, raw atmospheric seed and unknown high-order calibration map
are shared across cells. Spectrum branches use their fixed independent developer
gains; no test-parent or per-frame normalization occurs. Detector draws differ
by cell. All three flux replicas, both designs and four frozen methods give
160 condition-specific acquisitions,480 rows per design and3,840 finite estimates.
The source-matched native/GPU canary passed four image checks, maximum relative
$L_1=7.36\times10^{-17}$. Full diagnostics ended native exit zero in 281s wall.

The three-factor contrasts use $-1/+1$ coding. Reported effects are twice the
orthogonal regression coefficients; a main effect is the average high-minus-low
error across the other factors. Wavefront effects average the eight cell RMSEs,
and are distinct from the standardized-MSE contrasts. Parent bootstraps preserve
all eight cells and flux pairing; marginal 95\% intervals are exploratory and
unadjusted across methods, designs and main/interaction effects.

\begin{table}[ht]
\centering\small
\caption{Diagnostic wavefront-RMSE main effects in nm. Values show single/pair;
positive effects indicate degradation. Twenty independent paired parents.}
\begin{tabular}{@{}lrrr@{}}
\toprule
Method&Amplitude200$\to$300&Cutoff10$\to$6&Diversity1.0$\to$1.1\\
\midrule
MLP&27.68/24.74&16.36/12.01&$-0.62$/1.18\\
Compact CNN&27.04/25.73&17.10/10.61&$-0.35$/3.46\\
ResNet18&27.43/27.22&16.46/9.98&1.69/0.77\\
Ridge&17.94/15.19&9.80/7.06&$-2.74$/0.94\\
\bottomrule
\end{tabular}
\label{tab:factors}
\end{table}

Amplitude and spectrum effects are positive for every tested method/design
with exploratory intervals excluding zero. For MLP single images the intervals
are 22.19--33.02nm and10.89--21.56nm respectively; for pairs 20.10--29.72nm and
8.63--15.46nm. All neural diversity main-effect intervals include zero.
Several interactions remain visible, so main effects do not imply that diversity
is universally negligible or that effects simply add in every condition.
The unknown calibration map was held fixed rather than varied as a separate
factor; this study does not isolate calibration-error contributions or establish
a general physical mechanism outside the tested ranges.

The prospective continuation is to freeze validation-based method choices,
then test 100 untouched paired nominal/shifted parents in a separate seed
namespace. Four primary standardized-MSE contrasts would use 9,999 joint
parent-bootstrap draws and 98.75\% marginal intervals, an approximate Bonferroni
allocation conditional on fixed models/calibration. A bounded two-update
physical comparison would use the first 20 parents at 10,000 electrons. No such
confirmation or physical follow-on has run, and native admission remains required.

\section{Interpretation and evidence limits}
The present baseline supports a restricted result: a small trained estimator
has lower nominal point wavefront error and standardized MSE than the chosen
ridge control at sub-ms median inference cost. Physical fitting can further
reduce nominal error on the smaller subset at much greater computational cost.
A combined physical shift erodes the neural advantage over ridge on the full
population and worsens the tested physical-refinement errors. Increasing model
capacity improves nominal errors but has not resolved shifted behavior. The paired
diagnostics identify amplitude and spectrum changes as the largest average
effects within the specified ranges. Training that covers these nuisance
distributions is a motivated future experiment, not an established remedy.
Any training change requires independent parents, fixed calibration and
untouched evaluation; reusing inspected development outcomes cannot confirm it.

The simulations omit a measured telescope controller, telemetry-derived
residuals, scintillation, detector characterization and bench/sky reference
coefficients. The current experiment concerns four global modes; differential
segment/petal piston and fast tip--tilt--focus control require separate
measurement and dynamics studies. No current facility readiness, sky coverage,
closed-loop stability, exact prior-method superiority or publication novelty
claim follows from these results.

\section*{Reproducibility and drafting provenance}
\begin{sloppypar}
The local pipeline preserves source/configuration/runtime identities, basis
transforms, parent allocation, shard and model checksums, admission records,
latency reports, failures and cumulative-resource charges. The authoritative
baseline reports are \path|evaluation.json| and
\path|flux-evaluation.json| under the preserved local run
\path|keck-benchmark-05|. Dataset and model
release have not been performed. The current source uses Python 3.11.15,
HCIPy 0.7.1 and Torch 2.11.0+cu128 in the pinned project environment.
Codex assisted computational development, evidence integration and draft
preparation. Mojtaba's scientific review remains required; journal-specific
disclosure wording must be checked before any external submission.
\end{sloppypar}

\begin{thebibliography}{99}\small
\bibitem{fienup} J. R. Fienup (1982). Phase retrieval algorithms: a comparison.
\emph{Applied Optics} 21, 2758--2769. \href{https://doi.org/10.1364/AO.21.002758}{doi:10.1364/AO.21.002758}.
\bibitem{paxman} R. G. Paxman, T. J. Schulz and J. R. Fienup (1992). Joint estimation
of object and aberrations by using phase diversity. \emph{JOSA A} 9, 1072--1085.
\href{https://doi.org/10.1364/JOSAA.9.001072}{doi:10.1364/JOSAA.9.001072}.
\bibitem{meimon} S. Meimon, T. Fusco and L. M. Mugnier (2010). LIFT: a focal-plane
wavefront sensor for real-time low-order sensing on faint sources.
\emph{Optics Letters} 35, 3036--3038. \href{https://doi.org/10.1364/OL.35.003036}{doi:10.1364/OL.35.003036}.
\bibitem{plantet2011} C. Plantet, S. Meimon, T. Fusco and J. Conan (2011).
Experimental validation of the linearized focal-plane technique (LIFT).
\href{https://ao4elt2.lesia.obspm.fr/sites/ao4elt2/IMG/pdf/113meimon.pdf}{AO4ELT2 proceedings}.
\bibitem{plantet2013} C. Plantet, S. Meimon, J. Conan and T. Fusco (2013).
Experimental validation of LIFT for estimation of low-order modes in low-flux
wavefront sensing. \emph{Optics Express} 21, 16337--16352.
\href{https://doi.org/10.1364/OE.21.016337}{doi:10.1364/OE.21.016337}.
\bibitem{korkiakoski} V. Korkiakoski et al.\ (2014). Fast \& Furious focal-plane
wavefront sensing. \emph{Applied Optics} 53, 4565.
\href{https://arxiv.org/abs/1406.1006}{arXiv:1406.1006}.
\bibitem{bos} S. P. Bos et al.\ (2021). ``Fast'' and Furious focal-plane wavefront
sensing at W. M. Keck Observatory. Conference preprint.
\href{https://arxiv.org/abs/2107.07601}{arXiv:2107.07601}.
\bibitem{guyon} O. Guyon (2005). Limits of Adaptive Optics for high contrast imaging.
\emph{ApJ} 629, 592--614. \href{https://arxiv.org/abs/astro-ph/0505086}{arXiv:astro-ph/0505086}.
\bibitem{taheri} Mojtaba Taheri, M. Molahasani, S. Ragland, B. Neichel and
P. Wizinowich (2024). AI-Powered Low-Order Focal Plane Wavefront Sensing in
Infrared. \emph{SPIE} 13097, 1309783.
\href{https://arxiv.org/html/2410.12084v1}{arXiv:2410.12084}.
\bibitem{salgueiro} R. M. Salgueiro et al.\ (2026). Slow focus sensor for the Keck I
laser guide star adaptive optics system using focal plane wavefront sensing.
\href{https://arxiv.org/html/2602.15746v1}{arXiv:2602.15746}.
\bibitem{kuznetsov} A. Kuznetsov, S. Oberti, B. Neichel and T. Fusco (2024).
Striving towards robust phase diversity on-sky: Implementing LIFT for
VLT/MUSE-NFM. \emph{A\&A} 687, A221.
\href{https://arxiv.org/abs/2406.08529}{arXiv:2406.08529}.
\bibitem{orban} G. Orban de Xivry et al.\ (2021). Focal plane wavefront sensing
using machine learning: performance of convolutional neural networks compared
to fundamental limits. \emph{MNRAS} 505, 5702--5713.
\href{https://arxiv.org/abs/2106.04456}{arXiv:2106.04456}.
\bibitem{quesnel} M. Quesnel, G. Orban de Xivry, G. Louppe and O. Absil (2022).
A deep learning approach for focal-plane wavefront sensing using vortex phase
diversity. \href{https://arxiv.org/abs/2210.00632}{arXiv:2210.00632}.
\bibitem{wu} Y. Wu, Y. Guo, H. Bao and C. Rao (2020). Sub-Millisecond Phase
Retrieval for Phase-Diversity Wavefront Sensor. \emph{Sensors} 20, 4877.
\href{https://doi.org/10.3390/s20174877}{doi:10.3390/s20174877}.
\bibitem{por} E. H. Por, S. Y. Haffert, V. M. Radhakrishnan, D. S. Doelman,
M. van Kooten and S. P. Bos (2018). High Contrast Imaging for Python (HCIPy):
an open-source adaptive optics and coronagraph simulator. \emph{SPIE} 10703.
\href{https://ehpor.github.io/assets/pdf/Por-2018-HCIPy.pdf}{Author manuscript}.
\bibitem{he} K. He, X. Zhang, S. Ren and J. Sun (2016). Deep Residual Learning for
Image Recognition. \emph{CVPR}.
\href{https://arxiv.org/abs/1512.03385}{arXiv:1512.03385}.
\bibitem{adam} D. P. Kingma and J. Ba (2015). Adam: A Method for Stochastic
Optimization. \emph{ICLR}.
\href{https://arxiv.org/abs/1412.6980}{arXiv:1412.6980}.
\bibitem{bond} C. Z. Bond et al.\ (2020). Adaptive optics with an infrared
pyramid wavefront sensor at Keck. \emph{JATIS} 6, 039003.
\href{https://doi.org/10.1117/1.JATIS.6.3.039003}{doi:10.1117/1.JATIS.6.3.039003}.
\bibitem{landman} R. Landman et al.\ (2025). Making the unmodulated pyramid
wavefront sensor smart II. First on-sky demonstration of extreme adaptive
optics with deep learning. \emph{A\&A} 696, L1.
\href{https://arxiv.org/html/2503.16690v1}{arXiv:2503.16690}.
\bibitem{singh} G. Singh et al.\ (2026). Addressing the low-wind effect with the
Lyot-based low-order wavefront sensor on SCExAO. Preprint.
\href{https://arxiv.org/html/2609.00737v1}{arXiv:2609.00737}.
\end{thebibliography}
\end{document}
